\chapter{Several cuts: junctions, elimination, and refinement}\label{ch:10}

Put a second artificial cut into a regional problem. An observer outside the region should not be able to detect which cut was erased first. There are two tests of that expectation. The first asks whether the same points and fibres were identified. The second asks whether a signal launched with the same source and outer conditions gives the same response. A correct point-set quotient passes only the first test. A calculation on the wave equation below passes both and exhibits the travel time that would be lost by a careless partial solve.

\section{The geometry at a junction}

Consider finitely many oriented pieces with faces, common edges of faces, and possibly deeper intersections. A \emph{stratum} here means a smooth piece of fixed incidence type: one can move on it without changing which faces meet. For example the four quadrants of the spatial plane times a time interval give four timelike half-faces meeting on a timelike line. The line is present in the bookkeeping even when the unsplit spacetime is flat and has no object there.

Fix the two-sided face identifications, their lifts to primitive bundles, the identifications on the remaining strata, and the local sector to be assembled. Require the identifications on a common stratum to be consistent: two routes to the same local germ give the same identification of that germ and its fibre. This tests transition maps, not parallel transport by a generally curved connection. Require compatible collar and junction charts which assemble to the intended smooth local models, with the required matched jets of smooth data. Assume the quotient is Hausdorff, has the desired time orientation, and carries the stipulated physical boundary and causal problem. Arbitrarily identifying a set of half-faces need not give a manifold: a cross section normal to an ordinary interior junction must be a disk with the indicated sectors, not two disks joined only at their centres. Metric jets must match for a smooth metric even when the topology is that of a disk.

Let $X$ be the disjoint union of the pieces. Colour one group of identifications blue and the rest red. Two points have the same image in the simultaneous quotient precisely if they are joined by a finite chain of steps of either colour. In the blue quotient the blue steps become equalities; the red steps descend as a relation \emph{generated by their images}. Every chain in the iterated quotient lifts to red steps interspersed with blue equivalences, so defines exactly the same final equivalence relation. By the defining property of the quotient topology, a set in the result is open just when its preimage in $X$ is open, in either order. The resulting homeomorphism is smooth with smooth inverse in the fixed assembled charts, and respects the bundle transitions by their consistency. This proves order independence of this geometry. It would not justify smoothing each intermediate corner in a newly chosen way or forgetting a still exposed transition map.

A \emph{defect} is additional physics supported on a stratum: for example a prescribed jump law, source, or independent dynamical field. It is not another name for the stratum itself. To see the distinction at the level of equations, vary the regional actions retaining the remaining faces as open. Use outward response densities and express identified traces in a common fibre. Up to the bulk and external terms the sum takes the form
\begin{equation}
 \sum_F\int_F\langle\Pi_{F,1}+\Pi_{F,2},\delta q_F\rangle
 +\sum_J\int_J\langle k_J,\delta z_J\rangle+\delta S_{\rm def} .
 \label{eq:10-junction-first-variation}
\end{equation}
Here $F$ ranges over paired face interiors, $z_J$ collects the compatible traces at a junction $J$, and $k_J$ is the \emph{signed sum} of terms arriving there after any tangential integrations by parts, expressed in the same trace variables. If an action uses more traces or jets they must also be included. Vary in face interiors first to obtain the flux equations. At $J$ vary only the traces that extend as a common admissible variation to the incident faces. If they are constrained, it is the restriction of the covector in \eqref{eq:10-junction-first-variation} to that variation space that must vanish. Treating each copy of $\delta z_J$ as independent can create fictitious equations. For the first-derivative bulk scalar with no face action the first integration by parts gives only face fluxes, no $k_J$. For a common corner potential on adjacent ordinary faces, equal pullbacks with opposite incidences cancel (Appendix~\ref{app:A}). Neither example demands an extra force at an artificial junction.

By contrast, on a timelike junction with proper time $t$, an independent coordinate $z$ with action $\int_J(\dot z^2-\mu^2z^2)\,\mathrm{d}t/2$ contributes $-(\ddot z+\mu^2z)\delta z$ to the integrand and $[\dot z\,\delta z]_i^f$ at its ends. If a coupling $\int_J c z q_J\,\mathrm{d}t$ is well defined on the chosen trace space, it also contributes $c q_J\,\delta z+c z\,\delta q_J$. The oscillator equation and trace law have changed; the endpoint term changes the potential. Sewing the same geometry but dropping these terms is a different theory.

\section{Keep the other cut open}

The actual equations on two cuts may be nonlinear. Denote them by $F_1(x,y)=0$ and $F_2(x,y)=0$ for independent coordinates $x,y$ on the joint trace domain; external data are held fixed. If the first equation has exactly one solution $x=\chi(y)$ for each admissible $y$ under consideration, the remaining equation is $F_2(\chi(y),y)=0$. Its solutions correspond to simultaneous solutions by $y\mapsto(\chi(y),y)$, not by $y\mapsto(\chi(0),y)$. If $x$ is multivalued, the relation of allowed $(x,y)$ and its branches must be retained. This elementary observation explains why the intermediate object is an \emph{open family} over the unsolved history. It must still respond when that history is changed. In particular a solution produced by closing all the other faces with zero trace is not that family.

There is an important qualification hidden in the word ``independent''. At a common edge the two face histories might have to agree. When a continuous lifting of the common trace exists, one can parametrize each face by a lifted common trace plus a remainder with zero edge trace, and retain the shared variable exactly once. More generally work on the joint compatible trace space without a splitting. The following block proof applies when honest product coordinates have been supplied for the block being eliminated and the variables kept. The block kept can include junction variables and their equations. It does not apply to independently varying incompatible corner values.

For a linear (or linearized) causal problem use the construction of Chapter~\ref{ch:08} with all traces prescribed. A particular response $G f$ to the collection of sources $f$ leaves a mismatch $Nf$ at the faces. A homogeneous lift $H\eta$ of the unknown trace history changes the mismatch by $M\eta$. Thus the equation to be solved is $M\eta=-Nf$, not just equality of boundary point values. In coordinates $\eta=(x,y)$ write
\begin{equation}
 A x+B y=b,\qquad C x+D y=d,
 \qquad \binom{b}{d}=-Nf .
 \label{eq:10-blocks}
\end{equation}
For example $B$ is the response measured on the first face when the history on the second changes with the first held fixed. It can be nonlocal in time. The blocks map histories $X,Y$ to mismatch spaces $U,V$: $A:X\to U$, $B:Y\to U$, $C:X\to V$, $D:Y\to V$. These names stand for specified trace and response domains, including regularity, compatibility at initial corners, gauge conditions if used, and support. There is no presumption that a Dirichlet parametrization valid for a timelike face works on a null or a Cauchy face.

\begin{proposition}[Two-stage causal solve]
Use the retarded trace spaces, with zero history before the source, and suppose all the displayed block operations act continuously on their stated domains. Suppose $A$ has a two-sided inverse there and $S=D-CA^{-1}B:Y\to V$ does also. Assume the blocks and inverses used propagate supports only into their allowed causal futures, with domains stable under the compositions. Then the simultaneous system \eqref{eq:10-blocks} has a unique retarded solution given by
\begin{equation}
 y=S^{-1}(d-CA^{-1}b),\qquad x=A^{-1}b-A^{-1}B y .
 \label{eq:10-staged}
\end{equation}
If the reverse pivots $D$ and $A-BD^{-1}C$ meet the corresponding assumptions, reversal gives the same response, including on all external faces left open. The advanced statement uses causal pasts instead.
\end{proposition}
\begin{proof}
At fixed $y$ the first equation gives $x=A^{-1}(b-B y)$. Substitution in the second gives $C A^{-1}b+(D-CA^{-1}B)y=d$. This proves necessity of \eqref{eq:10-staged}; substitution back into both rows proves sufficiency. Zero right-hand side gives $y=0$ and then $x=0$. Thus any reverse solve satisfying the two rows gives this same pair. A retained external measurement of the form $m_0(f)+m_1 x+m_2 y$ also agrees; so does the reconstructed bulk field $G f+H(x,y)$. Each operation is causal by hypothesis, and the composition of allowed causal paths is causal in the assembled spacetime. A matching solution with this support gives the global retarded response when assembly regularity and uniqueness of the global problem hold. In the notation of a single interface solve $L$ this is precisely $G-HLN$, now computed by two substitutions.
\end{proof}

For finitely many blocks apply the argument repeatedly, preserving the updated inhomogeneous term as well as the operator on every remaining history. On a restricted physical forcing range one can weaken surjectivity, but then $b-B y$ must actually be in the range of $A$ for the retained $y$, and $d-CA^{-1}b$ in the range of $S$. A right inverse alone leaves possible homogeneous solutions; either retain them or prove independence of the fields and observables being compared. An algebraic inverse need not be causal: the inverse of a time delay on unrestricted histories is a time advance. Even on its range it need not preserve retarded support relative to the applied force.

\begin{remark}[A bad pivot and a real zero mode]
Let both unknowns be histories on one time axis. The time-local equations $y=b$, $x+y=d$ have the unique causal solution $(d-b,b)$, but the coefficient of $x$ in the first row is zero. Eliminating $x$ by inverting that coefficient is impossible, although eliminating $y$ using the second row works. This is a block counterexample, not a failure of the geometric quotient. A different issue occurs for $x-y=b$, $-x+y=d$. The range requires $b+d=0$, and the common shift of $x,y$ is undetermined. It is not legitimate to set that shift to zero merely to write an inverse: it might be a gauge redundancy, an unprescribed incoming datum, or a physical junction mode. The identification and its effect on observables require separate information.
\end{remark}

\section{A refinement that remembers the travel time}

Return to the line of Chapter~\ref{ch:08}, with equation $u_{tt}-u_{\xi\xi}=f$, wave speed one, zero past disturbances and outgoing solutions at both ends. Put cuts at positions $\xi=0$ and $\xi=\ell>0$, with histories $q(t),r(t)$ respectively. Take both history and mismatch spaces to be $\mathcal H_0=\{h\in C^\infty(\mathbb R):h(t)=0\ \text{for}\ t\leq0\}$, with convergence of every derivative on compact time intervals. Start the Laplace calculation with smooth compact driving functions; the resulting time-domain maps below extend to $\mathcal H_0$. Measure the mismatch by the sum of the \emph{outward spatial derivatives}. For the action with kinetic density $(u_t^2-u_\xi^2)/2$ the variational coefficient is the negative of this sum. Multiplying both the homogeneous mismatch and that of the particular field by $-1$ leaves the interface equation unchanged.

For zero interior source and prescribed transforms $\widehat q,\widehat r$, the transformed middle field at position $\xi$ is the elementary solution of $u''=s^2u$,
\[
 \widehat u(\xi)=\frac{\sinh(s(\ell-\xi))}{\sinh(s\ell)}\widehat q
              +\frac{\sinh(s\xi)}{\sinh(s\ell)}\widehat r,
 \qquad 0\leq\xi\leq\ell .
\]
The two exterior solutions are $e^{s\xi}\widehat q$ for $\xi\leq0$ and $e^{-s(\xi-\ell)}\widehat r$ for $\xi\geq\ell$. At zero the interior outward derivative is $-\widehat u'(0)=s\coth(s\ell)\widehat q-s\operatorname{csch}(s\ell)\widehat r$, and the left exterior supplies $+s\widehat q$. At $\ell$ the interior derivative is $+\widehat u'(\ell)$ and the right exterior supplies $+s\widehat r$. Thus the matrix carrying the two histories to the two summed derivative mismatches is
\[
 M(s)=\begin{pmatrix}A&B\\ B&A\end{pmatrix},\qquad
 A=s(1+\coth(s\ell)),\qquad B=-s\operatorname{csch}(s\ell).
\]
Inhomogeneous bulk fields with fixed traces are accounted for in the two right-hand sides $\widehat b,\widehat d$. In particular they are not included a second time in the homogeneous map $M$.

Three identities do the entire elimination:
\[
 A^{-1}=\frac{1-e^{-2s\ell}}{2s},\qquad A^{-1}B=-e^{-s\ell},
 \qquad A-BA^{-1}B=2s .
\]
No pivot vanishes for $\operatorname{Re}s>0$. More is true than invertibility of numbers at each frequency. Applied to a causal function $b$, the first inverse gives $\tfrac12\int_{t-2\ell}^t b(\tau)\,\mathrm{d}\tau$ with zero extension to negative times. The factor $e^{-s\ell}$ gives a delay by $\ell$; $1/(2s)$ is half the integral from the past. The direct blocks also act causally: expanding $(1-e^{-2s\ell})^{-1}$ as a geometric series gives delayed derivatives with only finitely many arrivals on each bounded time interval. Thus all these maps are continuous on $\mathcal H_0$ and the inverse identities extend there: on any finite interval cut off the data smoothly beyond its end, where no change can affect the answer. The full solution, displayed in transform notation on data for which the transform exists, is
\begin{equation}
 \binom{\widehat q}{\widehat r}
 =\frac{1}{2s}
  \begin{pmatrix}1&e^{-s\ell}\\ e^{-s\ell}&1\end{pmatrix}
  \binom{\widehat b}{\widehat d} .
 \label{eq:10-two-cuts-line}
\end{equation}

Now view zero as a subdivision of the left half of a problem originally cut only at $\ell$. Eliminate $q$ while $r$ is still free. The surviving equation is $2s\widehat r=\widehat d+e^{-s\ell}\widehat b$. This is exactly the old one-cut operator: each of the two unbroken half-lines at $\ell$ contributes $s\widehat r$. A mismatch at the new face contributes to the remaining source only after time $\ell$.

Does the \emph{forcing} produced by a regional bulk experiment also reassemble correctly? Place an arbitrary smooth compact source $f(t,\xi)$ inside $0<\xi<\ell$ after time zero. Call the two coefficients of $\widehat q,\widehat r$ in the displayed middle field $h_0(\xi),h_\ell(\xi)$. The zero-trace particular solution $v$ in the middle satisfies $(s^2-\partial_\xi^2)v=\widehat f$. Integration by parts gives $\int_0^\ell h\widehat f\,\mathrm{d}\xi=[-h v'+h'v]_0^\ell$ for either coefficient. Thus $\widehat b=v'(0)=\int h_0\widehat f$ and $\widehat d=-v'(\ell)=\int h_\ell\widehat f$ are precisely the negatives of its outward derivatives. The combination $h_\ell+e^{-s\ell}h_0=e^{-s(\ell-\xi)}$ follows by substitution, or by solving the homogeneous equation with these two endpoint values. It gives the measured field at the retained cut:
\[
 \widehat r=\int_0^\ell \frac{e^{-s(\ell-\xi)}}{2s}\,\widehat f(s,\xi)\,\mathrm{d}\xi .
\]
This is the uncut answer. Indeed the full-line Green function for $s^2-\partial_\xi^2$ is $e^{-s|\xi-\xi'|}/(2s)$: it decays at both ends, is continuous, and its first derivative jumps by $-1$ at $\xi'$. In time, a signal from $\xi$ cannot arrive before $\ell-\xi$. Checking only that the homogeneous operator returned to $2s$ would not have tested this source dependence.

\begin{exercise}[Erase the other cut]
For the wave example retain $q$ instead of $r$. Write both the remaining source and the open lift of $r$ at arbitrary $q$. \emph{Answer:} $2s\widehat q=\widehat b+e^{-s\ell}\widehat d$ and $\widehat r=A^{-1}\widehat d+e^{-s\ell}\widehat q$.
\end{exercise}

\begin{exercise}[A branch lost by partial solving]
At a junction suppose the two equations are $x^2=y$ and $x+y=0$ for real traces. Compare the simultaneous solutions with a procedure retaining only $x=\sqrt y$. \emph{Hint:} solve $x^2+x=0$. The negative root is not a second description of the positive one; it gives a different matched trace.
\end{exercise}
